% ==============================================================================
% Kapitel 4: Integration von Lexer- und Parser-Werkzeugen
% ==============================================================================
\chapter{Integration von Lexer- und Parser-Werkzeugen}
\label{chap:toolchains}

Der größte Teil von Stroetmanns Skript, die Kapitel~6 bis~11, ist der Syntaxanalyse kontextfreier Sprachen gewidmet~\cite{stroetmann2026formal}.
Das Skript verbindet dabei zwei Zugänge: Die Studierenden benutzen einen Parsergenerator als Werkzeug, um einen Taschenrechner, ein Programm zur symbolischen Differentiation, einen Interpreter und schließlich den Integer-C-Compiler zu bauen, und sie implementieren die zugrunde liegenden Parsing-Verfahren selbst~\cite{earley1970efficient,knuth1965translation}.

Als Parsergenerator verwendet der Kurs inzwischen \emph{Larkita}, eine von Stroetmann entwickelte, schlanke Fassung der Bibliothek Lark.
Die Notizbücher der Vorlesung bilden beide Zugänge ab:
\begin{enumerate}
  \item \textbf{Generierte Parser:} Wie binden sich Larkita in Python und seine Portierungen nach OCaml und F\# in interaktive Notizbücher ein, wie entsteht aus einer Grammatik ohne semantische Aktionen ein abstrakter Syntaxbaum, wie wird dieser in den drei Typsystemen dargestellt, und wie unterstützen die Werkzeuge bei Grammatikkonflikten?
  \item \textbf{Didaktische Eigenimplementierungen:} Wie werden Parsing-Verfahren von Grund auf im Quelltext realisiert: vom rekursiven Abstieg über den Earley-Algorithmus für beliebige kontextfreie Grammatiken bis zur Erzeugung von SLR- und kanonischen LR-Tabellen für eine Grammatik der Sprache C?
\end{enumerate}

In diesem Kapitel vergleichen wir beide Ebenen über alle drei Sprachwelten hinweg.
Frühere Fassungen der Notizbücher verwendeten Lark in Python, \texttt{ocamllex} und Menhir in OCaml sowie \texttt{FsLex} und \texttt{FsYacc} in F\#; diese Werkzeuge dienen im Folgenden als Vergleichsmaßstab.

% ------------------------------------------------------------------------------
\section{Larkita: Grammatiken ohne semantische Aktionen}
\label{sec:larkita}
% ------------------------------------------------------------------------------

Larkita unterscheidet sich in drei Punkten von Lark: Es erzeugt stets einen \emph{kanonischen LR(1)-Parser} (auf Wunsch mit einer sicheren Verschmelzung von Zuständen mit gleichem Kern), der Parser liefert den abstrakten Syntaxbaum als \emph{verschachteltes Tupel}, und jede Erzeugung schreibt eine \emph{Ausgabedatei} mit allen LR(1)-Zuständen samt Vorausschaumengen.
Eine Grammatik enthält Token und Regeln, aber keine semantischen Aktionen; die Form des Syntaxbaums folgt aus der Grammatik selbst:
Die Zeichenkettenliterale einer Alternative bilden den \emph{Kopf} des Tupels, alle anderen Symbole seine \emph{Kinder}.
So liefert \texttt{expr "+" product} das Tupel \texttt{('+', ast(expr), ast(product))}, die Zuweisung \texttt{IDENTIFIER ":=" expr ";"} das Tupel \texttt{((':=', ';'), ('IDENTIFIER', x), ast(expr))} und ein benanntes Token das Paar \texttt{('NUMBER', '42')}.
Eine Alternative aus genau einem Symbol reicht dessen Baum durch, \texttt{\%ignore "(", ")"} entfernt Literale aus dem Kopf, und ein Alias \texttt{-> name} ersetzt ihn.

Für die OCaml- und die F\#-Fassung des Kurses wurde Larkita vollständig portiert (\texttt{ocaml/lib/larkita*.ml}, \texttt{f\#/Support/Larkita}).
Beide Portierungen lesen dieselben \texttt{.lark}-Dateien wie die Python-Notizbücher, berechnen denselben kanonischen LR(1)-Automaten und schreiben eine bis aufs Byte identische Ausgabedatei.
Die Terminale werden mit PCRE2 (OCaml) bzw.\ mit \texttt{System.Text.RegularExpressions} (F\#) erkannt, deren Syntax für reguläre Ausdrücke der von Python weitgehend entspricht.
Ein Differenztest zeichnet das Verhalten der Python-Bibliothek auf (ihre Einheitstests, die Beispiele des Larkita-Tutorials und alle Kursnotizbücher) und spielt jeden Aufruf mit den Portierungen nach: Ausgabedateien, Syntaxbäume, Token und Fehlermeldungen stimmen für alle 282 Grammatiken überein.

\begin{table}[htbp]
  \centering
  \small
  \caption{Architektonischer Vergleich der Parser-Werkzeugketten im Curriculum.}
  \label{tab:toolchain_architectures}
  \begin{tabularx}{\textwidth}{lXXX}
    \toprule
    \textbf{Dimension} & \textbf{Python (Larkita)} & \textbf{OCaml (Larkita-Portierung)} & \textbf{F\# (Larkita-Portierung)} \\
    \midrule
    \textbf{Architekturtyp} & Bibliothek, Generierung zur Laufzeit & Bibliothek im Toplevel, Generierung zur Laufzeit & Skriptbibliothek in F\# Interactive, Generierung zur Laufzeit \\
    \addlinespace[3pt]
    \textbf{Spezifikation} & \texttt{.lark}-Datei (EBNF) ohne semantische Aktionen & dieselbe \texttt{.lark}-Datei & dieselbe \texttt{.lark}-Datei \\
    \addlinespace[3pt]
    \textbf{Verfahren} & Kanonisches LR(1), optional sichere Zustandsverschmelzung & identisch & identisch \\
    \addlinespace[3pt]
    \textbf{Lexer} & Python \texttt{re}, kontextabhängig & PCRE2, kontextabhängig & .NET \texttt{Regex}, kontextabhängig \\
    \addlinespace[3pt]
    \textbf{Syntaxbaum} & Verschachtelte Tupel ohne Typprüfung & Summentyp \texttt{ast}, Umwandlung in getypte ADTs & Union \texttt{Ast}, Umwandlung in getypte ADTs \\
    \addlinespace[3pt]
    \textbf{Konfliktdiagnose} & Ausgabedatei \texttt{.out}; \texttt{ConflictWarning}, \texttt{GrammarError} & identische Ausgabedatei; Warnung bzw.\ \texttt{Conflict\_warning}, \texttt{Grammar\_error} & identische Ausgabedatei; Warnung bzw.\ \texttt{ConflictWarning}, \texttt{GrammarError} \\
    \addlinespace[3pt]
    \textbf{Interaktivität} & Sofortige Zellauswertung & Sofortige Zellauswertung & Sofortige Zellauswertung \\
    \bottomrule
  \end{tabularx}
\end{table}

\begin{figure}[htbp]
  \centering
  \begin{tikzpicture}[
    node distance=1.1cm and 1.15cm,
    box/.style={rectangle, draw=black!70, rounded corners=3pt, fill=CodeBg, inner sep=5pt, font=\scriptsize, align=center},
    pybox/.style={box, fill=PythonBlue!10, draw=PythonBlue!80},
    ocbox/.style={box, fill=OCamlOrange!10, draw=OCamlOrange!80},
    fsbox/.style={box, fill=FSharpPurple!10, draw=FSharpPurple!80},
    arrow/.style={-{Stealth[scale=0.9]}, draw=black!70, thick}
  ]
    % Python pipeline
    \node[pybox] (p1) {\textbf{Grammatikdatei}\\\texttt{calc.lark} (EBNF)};
    \node[pybox, right=of p1] (p2) {\textbf{Larkita}\\kanonische LR(1)-Tabelle\\und \texttt{calc.out}};
    \node[pybox, right=of p2] (p3) {\textbf{Syntaxbaum}\\verschachtelte Tupel\\ohne Typprüfung};
    \draw[arrow] (p1) -- node[above, font=\tiny] {\texttt{Larkita.open}} (p2);
    \draw[arrow] (p2) -- node[above, font=\tiny] {\texttt{parse}} (p3);

    % OCaml pipeline
    \node[ocbox, below=1.2cm of p1] (o1) {\textbf{Grammatikdatei}\\dieselbe \texttt{calc.lark}};
    \node[ocbox, right=of o1] (o2) {\textbf{Larkita-Portierung}\\reines OCaml, PCRE2\\identische \texttt{calc.out}};
    \node[ocbox, right=of o2] (o3) {\textbf{Syntaxbaum}\\Summentyp \texttt{ast}\\$\to$ getypte ADT};
    \draw[arrow] (o1) -- node[above, font=\tiny] {\texttt{Larkita.open\_}} (o2);
    \draw[arrow] (o2) -- node[above, font=\tiny] {\texttt{Larkita.parse}} (o3);

    % F# pipeline
    \node[fsbox, below=1.2cm of o1] (f1) {\textbf{Grammatikdatei}\\dieselbe \texttt{calc.lark}};
    \node[fsbox, right=of f1] (f2) {\textbf{Larkita-Portierung}\\F\#, .NET \texttt{Regex}\\identische \texttt{calc.out}};
    \node[fsbox, right=of f2] (f3) {\textbf{Syntaxbaum}\\Union \texttt{Ast}\\$\to$ getypte ADT};
    \draw[arrow] (f1) -- node[above, font=\tiny] {\texttt{Larkita.Open}} (f2);
    \draw[arrow] (f2) -- node[above, font=\tiny] {\texttt{parse}} (f3);
  \end{tikzpicture}
  \caption{Erzeugung und Anwendung der Parser in Python, OCaml und F\#: alle drei Fassungen lesen dieselbe Grammatikdatei.}
  \label{fig:toolchain_pipelines}
\end{figure}

Mit Larkita verschwinden die architektonischen Unterschiede, die die früheren Werkzeugketten prägten (\Cref{tab:toolchain_architectures} und \Cref{fig:toolchain_pipelines}).
Lark erzeugte den Parser zur Laufzeit in der Notizbuch-Zelle; OCaml schrieb getrennte Spezifikationen für \texttt{ocamllex} und Menhir in Dateien, rief beide Werkzeuge als Subprozesse auf und lud die erzeugten Module in den Toplevel; in F\# lagen die \texttt{.fsl}- und \texttt{.fsy}-Dateien in eigenen Projekten, die vor dem Start des Notizbuchs mit MSBuild übersetzt werden mussten.
Nun wird der Parser in allen drei Sprachen in derselben Zelle und aus derselben Grammatikdatei erzeugt, und eine Änderung an der Grammatik wird beim nächsten Zellendurchlauf wirksam.
Die semantischen Aktionen, die in Menhir und FsYacc den Parse-Baum bzw.\ den getypten Syntaxbaum aufbauten, entfallen ebenso wie die Token-Deklarationen; die Lexer- und Parser-Spezifikationen der Notizbücher schrumpfen auf die Grammatik des Skripts.

% ------------------------------------------------------------------------------
\section{Darstellung der abstrakten Syntaxbäume}
\label{sec:ast_representation}
% ------------------------------------------------------------------------------

Die verschachtelten Tupel von Larkita sind \emph{heterogen}: Ihr Kopf ist eine Zeichenkette oder ein Tupel von Zeichenketten, ihre Kinder sind Tupel, Token-Paare oder einzelne Literale, und ihre Länge hängt von der Alternative ab.
Python verarbeitet sie direkt mit strukturellem Mustervergleich, etwa \texttt{case ((':=', ';'), ('IDENTIFIER', v), e)}.
OCaml und F\# kennen keine Tupel variabler Länge; beide Portierungen stellen die Bäume deshalb durch einen Summentyp dar:
\lstinputlisting[style=notebook,language=Caml,title={\textbf{OCaml} -- \texttt{lib/larkita\_base.ml} (Auszug)}]{listings/04-larkita_ast-ocaml.ml}
Das Muster der Zuweisung lautet dann \texttt{Node ([":="; ";"], [Tok ("IDENTIFIER", v); e])}; ein Drucker zeigt die Werte in der Python-Notation an, sodass die Ausgaben der Notizbücher übereinstimmen.
Diese Darstellung ist ebenso ungetypt wie die Python-Tupel: Der Typ \texttt{ast} schränkt die Form eines Baums nicht ein, und eine Auswertungsfunktion braucht wie in Python einen letzten Fall für unbekannte Strukturen.
Die Taschenrechner und Interpreter der OCaml-Fassung folgen dennoch direkt den Python-Notizbüchern.
\Cref{lst:calculator_eval} zeigt die Auswertung des Taschenrechners aus \texttt{Chapter-07/01-Calculator.ipynb} in allen drei Sprachen; die Muster entsprechen einander Fall für Fall.

\begin{figure}[htbp]
  \centering
\lstinputlisting[style=notebook,language=Python,title={\textbf{Python} -- \texttt{Chapter-07/01-Calculator.ipynb}}]{listings/04-calculator_eval-python.py}
\lstinputlisting[style=notebook,language=Caml,title={\textbf{OCaml} -- \texttt{Chapter-07/01-Calculator.ipynb}}]{listings/04-calculator_eval-ocaml.ml}
\lstinputlisting[style=notebook,language=FSharp,title={\textbf{F\#} -- \texttt{Chapter-07/01-Calculator.ipynb}}]{listings/04-calculator_eval-fsharp.fs}
  \caption{Auswertung der von Larkita gelieferten Syntaxbäume im Taschenrechner.}
  \label{lst:calculator_eval}
\end{figure}

Wo sich die statische Typprüfung lohnt -- beim Differenzierer, bei den Interpretern der F\#-Fassung und bei den Compilern --, wird der Baum einmal in getypte algebraische Datentypen umgewandelt; diese Umwandlung ersetzt die semantischen Aktionen der früheren Werkzeugketten, und danach prüft der Compiler die Vollständigkeit jedes Mustervergleichs.

% ------------------------------------------------------------------------------
\section{Konfliktdiagnose und Zustandsverschmelzung}
\label{sec:conflict_diagnostics}
% ------------------------------------------------------------------------------

Stroetmanns Skript widmet den Schiebe-Reduziere- und Reduziere-Reduziere-Konflikten große Teile von Kapitel~10 und~11~\cite{stroetmann2026formal}; die Notizbücher im Verzeichnis \texttt{Chapter-11} (in F\# \texttt{Chapter-10}) führen die Beispiele vor.
Da alle drei Fassungen denselben Generator verwenden, ist auch die Diagnose dieselbe.
Larkita löst jeden Schiebe-Reduziere-Konflikt zugunsten des Schiebens auf und meldet ihn mit einer Warnung; ein Reduziere-Reduziere-Konflikt, den keine Regelpriorität auflöst, ist ein Fehler.
Die Ausgabedatei listet die Regeln mit dem jeweils erzeugten Syntaxbaum, alle Konflikte samt der verworfenen Alternativen und alle Zustände mit ihren markierten Regeln und Vorausschaumengen.
Python macht die Warnung mit \texttt{warnings.simplefilter('error', ConflictWarning)} zu einer Ausnahme; die Portierungen bieten dafür die Funktion \texttt{with\_conflict\_warnings} bzw.\ \texttt{withConflictWarnings}.

Die früheren Werkzeuge unterschieden sich hier deutlich:
Lark löste die Konflikte seines LALR-Parsers stillschweigend auf und meldete sie nur über den Logger \texttt{lark.logger}; Menhir erklärte jeden Konflikt mit \texttt{--explain} anhand von Ableitungsbeispielen und zeigte in seinen Automaten nur die Kernelemente; FsYacc schrieb die Zustände in eine Protokolldatei.
Die Nummern der Zustände waren deshalb zwischen den Sprachen nicht vergleichbar; mit Larkita stimmen sie überein.

\paragraph{Der mysteriöse Reduziere-Reduziere-Konflikt.}
Ein didaktisch besonders interessantes Beispiel aus Stroetmanns Skript (Abschnitt~11.3.2, Notizbuch \texttt{05-Mysterious-Conflicts.ipynb}) ist der sogenannte mysteriöse Reduziere-Reduziere-Konflikt.
Die Grammatik ist eine LR(1)-Grammatik: In den beiden LR(1)-Zuständen $\{a \to X\bullet : \texttt{'y'},\ b \to X\bullet : \texttt{'z'}\}$ und $\{b \to X\bullet : \texttt{'y'},\ a \to X\bullet : \texttt{'z'}\}$ sind die Vorausschaumengen disjunkt.
Beide Zustände haben jedoch denselben Kern und werden bei der Konstruktion der LALR-Tabelle verschmolzen; im verschmolzenen Zustand kann der Parser bei \texttt{'y'} und \texttt{'z'} nicht mehr zwischen den beiden Reduktionen entscheiden.
Wie das Skript betont, lässt sich ein solcher Konflikt nicht durch Umschreiben der Grammatik beheben, weil der Fehler nicht in der Grammatik, sondern in der Verschmelzung der Zustände liegt.

Larkita berechnet den kanonischen LR(1)-Automaten; der Konflikt tritt deshalb in keiner der drei Fassungen auf.
Mit \texttt{merge\_states=True} (\texttt{\textasciitilde merge\_states:true}, \texttt{mergeStates = true}) verschmilzt Larkita Zustände mit gleichem Kern nur, wenn dadurch kein neuer Konflikt entsteht und der kontextabhängige Lexer dieselben Token liefert.
Das neue Notizbuch \texttt{06-Merging-States.ipynb} zeigt beides: Für die geschichtete Ausdrucksgrammatik schrumpft der Automat von 36 auf 19 Zustände, während die beiden kritischen Zustände der mysteriösen Grammatik getrennt bleiben und die Ausgabedatei den Grund dafür angibt.
Zum Vergleich: Lark brach mit einem \texttt{GrammarError} ab, Menhir lehnte die Grammatik mit \texttt{--lalr --strict} ab und baute sie mit \texttt{--canonical}, und FsYacc meldete zwei Reduziere-Reduziere-Konflikte, baute den Parser trotzdem und verlor die Eingaben \texttt{vxz} und \texttt{wxy}.

\paragraph{Präzedenzen.}
Larkita unterstützt wie Lark keine Präzedenzdeklarationen; Präzedenz und Assoziativität werden deshalb in allen drei Fassungen wie im Skript durch eine Schichtung der Grammatik mit einer Variablen je Präzedenzstufe kodiert.
Menhir, FsYacc und PLY besitzen dagegen yacc-artige Präzedenzdeklarationen, die die früheren Fassungen der Compiler-Notizbücher nutzten.
Das Dangling-Else-Problem des Integer-C-Compilers bleibt der einzige Konflikt der Compiler-Grammatik; Larkita löst ihn durch Schieben so auf, dass ein \texttt{else} zum innersten \texttt{if} gehört.

% ------------------------------------------------------------------------------
\section{Didaktische Eigenimplementierungen von Parsern}
\label{sec:handcrafted_parsers}
% ------------------------------------------------------------------------------

Neben der Nutzung vorgefertigter Generatoren implementieren die Studierenden im Verlauf der Vorlesung drei eigenständige Parsing-Engines:

\subsection{Rekursiver Abstieg}
Kapitel~6 des Skripts entwickelt einen Top-down-Parser mit rekursivem Abstieg für arithmetische Ausdrücke~\cite{stroetmann2026formal}.
Weil dieses Verfahren für linksrekursive Grammatiken nicht funktioniert, ersetzt Stroetmann die Linksrekursion durch die Iteration einer EBNF-Grammatik.
Bereits seine Python-Fassung formuliert den Parser als gegenseitig rekursive Funktionen, die jeweils ein Paar aus Ergebnis und verbleibender Tokenliste zurückgeben und die Fälle über Listenmuster unterscheiden (\texttt{case '+', *RL}).
Die F\#-Fassung übernimmt diese Struktur direkt als Gruppe \texttt{let rec \dots{} and \dots}, während die OCaml-Fassung wie bei der Thompson-Konstruktion (\cref{sec:case_thompson}) Vorwärtsreferenzen anlegt, damit die Funktionen in derselben Reihenfolge wie im Python-Notizbuch erklärt werden können.

\subsection{Allgemeines kontextfreies Parsing nach Earley}
Um beliebige kontextfreie Grammatiken---auch mehrdeutige, linksrekursive oder solche mit Epsilon-Regeln---verarbeiten zu können, behandelt Kapitel~9 des Skripts den Earley-Algorithmus~\cite{earley1970efficient}; die Python-Notizbücher dazu liegen im Verzeichnis \texttt{Chapter-09}.
Ein Earley-Objekt $\langle a \to \beta \bullet \gamma, k \rangle$ besteht aus einer Grammatikregel mit Markierungspunkt für die aktuelle Leseposition und einem Index für die Startposition der Regel im Eingabetext.
Für jede Position $j$ der Eingabe verwaltet der Algorithmus eine Menge $Q_j$ von Earley-Objekten, die das Skript mit drei Operationen aufbaut: \emph{Reading} (Einlesen des nächsten Tokens), \emph{Prediction} (Vorhersage der Regeln für eine erwartete Variable) und \emph{Completion} (Fortsetzen übergeordneter Regeln, sobald eine Regel vollständig erkannt ist).
Laut Skript hat der Algorithmus im allgemeinen Fall eine Komplexität von $O(n^3)$ und für eindeutige Grammatiken von $O(n^2)$.

\subsection{Generatoren für SLR(1)- und LR(1)-Tabellen}
In Kapitel~10 des Skripts konstruieren die Studierenden eigene Tabellengeneratoren für SLR- und kanonische LR-Parser (Notizbücher im Verzeichnis \texttt{Chapter-10}): Sie berechnen die Hülle von Mengen markierter Regeln, bestimmen First- und Follow-Mengen über Fixpunktberechnungen und füllen die Aktions- und Sprungtabellen.
Als Härtetest dient eine Grammatik für die Sprache C mit 63 Nichtterminalen und 211 Regelalternativen.
Laut Skript kommt ein SLR-Parser für diese Grammatik mit 349 Zuständen aus, erzeugt aber Konflikte, weil C keine SLR-Sprache ist, während der kanonische LR-Parser 1572 Zustände benötigt.
Für die SLR-Tabelle verwendet Stroetmann deshalb eine angepasste Grammatik ohne Typumwandlungen, Sprungmarken und \texttt{if}-Anweisungen ohne \texttt{else}-Zweig.
Das Dangling-Else-Problem behandelt das Skript anhand der Grammatik von Integer-C.


% ------------------------------------------------------------------------------
\section{Empirische Performanz der Parser-Werkzeuge}
\label{sec:toolchain_empirical}
% ------------------------------------------------------------------------------

Um die Rechenleistung zu vergleichen, wurden mit der Taschenrechnergrammatik \texttt{calc.lark} des Kurses (11 Produktionen) synthetische Summenausdrücke der Form \texttt{1 + 1 + \dots;} mit bis zu 10.000 Token analysiert (\Cref{tab:parser_throughput} und \Cref{fig:parser_throughput_scaling}; Skript \texttt{benchmarks/larkita\_benchmarks.py}).
Gemessen wurden die Erzeugung des kanonischen LR(1)-Parsers einschließlich Ausgabedatei und das Parsen einschließlich Lexer und Aufbau des abstrakten Syntaxbaums.
Jeder Wert ist der Mittelwert über 20 Durchläufe nach einem Aufwärmlauf.
Die OCaml-Portierung wurde zweimal gemessen: nativ kompiliert, vergleichbar mit der früheren Menhir-Messung, und im Bytecode-Toplevel, in dem auch die Notizbücher laufen.
Die F\#-Portierung lief in F\# Interactive wie die F\#-Notizbücher.
Zum Vergleich enthält die Tabelle die Messungen der früheren Werkzeugketten mit derselben Eingabe.

\begin{table}[htbp]
  \centering
  \footnotesize
  \setlength{\tabcolsep}{3pt}
  \caption{Gemessene Durchsatzwerte (Token pro Sekunde) und Erzeugungszeiten.
  Oben Larkita und seine Portierungen, unten die früheren Werkzeugketten als Referenz (Python 3.13, OCaml 5.4.1, .NET 9).}
  \label{tab:parser_throughput}
  \begin{tabularx}{\textwidth}{Xrrrrrrr}
    \toprule
    & \multicolumn{6}{c}{\textbf{Durchsatz in kTok/s bei \dots{} Token}} & \textbf{Erzeugung} \\
    \cmidrule(lr){2-7}
    \textbf{Werkzeug} & \textbf{100} & \textbf{500} & \textbf{1\,000} & \textbf{2\,500} & \textbf{5\,000} & \textbf{10\,000} & \textbf{ms} \\
    \midrule
    Larkita (Python)                  & 616,9 & 623,8 & 635,3 & 632,2 & 637,8 & 621,8 & 3,66 \\
    Larkita-Portierung OCaml, nativ   & \textbf{1\,186,2} & \textbf{1\,217,4} & \textbf{1\,324,5} & \textbf{1\,521,7} & \textbf{1\,675,6} & \textbf{1\,674,5} & \textbf{1,45} \\
    Larkita-Portierung OCaml, Toplevel & 529,9 & 538,6 & 543,1 & 536,4 & 521,9 & 533,1 & 3,55 \\
    Larkita-Portierung F\#, F\# Interactive & 492,6 & 517,0 & 514,4 & 539,3 & 476,8 & 627,5 & 2,96 \\
    \midrule
    Lark LALR (Python)                & 409,5 & 409,0 & 388,5 & 384,2 & 396,5 & 378,0 & 3,67 \\
    ocamllex + Menhir, nativ          & 31\,300,8 & 31\,207,6 & 30\,305,7 & 27\,349,4 & 33\,554,4 & 27\,456,8 & 14,63 \\
    FsLex + FsYacc (CoreCLR)          & 1\,815,5 & 515,6 & 3\,340,9 & 2\,868,3 & 3\,735,4 & 3\,562,2 & 1\,491,81 \\
    \bottomrule
  \end{tabularx}
\end{table}

\begin{figure}[htbp]
  \centering
  \begin{tikzpicture}
    \begin{axis}[
      width=0.92\textwidth,
      height=7.2cm,
      xmode=log,
      ymode=log,
      log basis x=10,
      log basis y=10,
      xlabel={Anzahl der Eingabetoken (logarithmisch)},
      ylabel={Parsing-Durchsatz in Token/Sekunde (logarithmisch)},
      xmin=80, xmax=12000,
      ymin=200000, ymax=60000000,
      grid=both,
      grid style={dashed,gray!35},
      legend style={at={(0.03,0.97)}, anchor=north west, font=\scriptsize},
      thick
    ]
      \addplot+[mark=*, color=OCamlOrange!90!black, line width=1.3pt] coordinates {
        (100, 1186240) (500, 1217434) (1000, 1324503) (2500, 1521699) (5000, 1675603) (10000, 1674481)
      };
      \addlegendentry{OCaml 5.4 (Larkita-Portierung, nativ)}

      \addplot+[mark=o, color=OCamlOrange!60, line width=1.0pt] coordinates {
        (100, 529942) (500, 538561) (1000, 543124) (2500, 536354) (5000, 521915) (10000, 533095)
      };
      \addlegendentry{OCaml 5.4 (Larkita-Portierung, Toplevel)}

      \addplot+[mark=triangle*, color=FSharpPurple, line width=1.3pt] coordinates {
        (100, 492611) (500, 517010) (1000, 514377) (2500, 539328) (5000, 476799) (10000, 627530)
      };
      \addlegendentry{F\# 9.0 (Larkita-Portierung, F\# Interactive)}

      \addplot+[mark=square*, color=PythonBlue, line width=1.3pt] coordinates {
        (100, 616903) (500, 623830) (1000, 635324) (2500, 632207) (5000, 637820) (10000, 621798)
      };
      \addlegendentry{Python 3.13 (Larkita)}

      \addplot+[mark=none, color=OCamlOrange!90!black, dashed, line width=0.9pt] coordinates {
        (100, 31300776) (500, 31207619) (1000, 30305665) (2500, 27349400) (5000, 33554432) (10000, 27456821)
      };
      \addlegendentry{Referenz: ocamllex + Menhir, nativ}

      \addplot+[mark=none, color=FSharpPurple, dashed, line width=0.9pt] coordinates {
        (100, 1815541) (500, 515592) (1000, 3340906) (2500, 2868288) (5000, 3735357) (10000, 3562167)
      };
      \addlegendentry{Referenz: FsLex + FsYacc}

      \addplot+[mark=none, color=PythonBlue, dashed, line width=0.9pt] coordinates {
        (100, 409472) (500, 409046) (1000, 388510) (2500, 384194) (5000, 396454) (10000, 377977)
      };
      \addlegendentry{Referenz: Lark LALR}
    \end{axis}
  \end{tikzpicture}
  \caption{Skalierungsverhalten des Parsing-Durchsatzes über steigende Eingabegrößen; gestrichelt die früheren Werkzeugketten.}
  \label{fig:parser_throughput_scaling}
\end{figure}

Die Ergebnisse zeigen zweierlei.
Erstens liegen Larkita und seine Portierungen in der Konfiguration der Notizbücher eng beieinander: Python erreicht rund 620.000 Token pro Sekunde, die OCaml-Portierung im Bytecode-Toplevel rund 530.000 und die F\#-Portierung in F\# Interactive 480.000 bis 630.000.
Larkita ist damit etwa anderthalbmal so schnell wie Lark, dessen LALR-Parser einen generischen Parse-Baum aufbaute.
Nativ kompiliert erreicht die OCaml-Portierung 1,2 bis 1,7 Millionen Token pro Sekunde, und sie erzeugt den Parser mit 1,45 Millisekunden am schnellsten.
Zweitens bleibt der Abstand zu den Generatoren, die Quelltext erzeugen, beträchtlich: Menhir übersetzt den Automaten in kompakten OCaml-Code und ist nativ etwa zwanzigmal schneller als die native Larkita-Portierung, und FsYacc ist etwa sechsmal schneller als die F\#-Portierung.
Die Portierungen wurden auf Übereinstimmung mit Python optimiert, nicht auf Geschwindigkeit; sie interpretieren die Parse-Tabelle zur Laufzeit, die OCaml-Portierung hält ihren Stapel in unveränderlichen Listen, und der Lexer wendet für jedes Token einen regulären Ausdruck an, der alle im aktuellen Zustand erlaubten Terminale als Alternativen enthält.
Für die kurzen Beispiele der Vorlesung -- selbst die Programme des Integer-C-Compilers umfassen nur einige hundert Token -- ist dieser Abstand ohne Bedeutung, während die Erzeugung des Parsers in allen drei Sprachen in wenigen Millisekunden geschieht statt in anderthalb Sekunden MSBuild-Lauf wie bei FsYacc.

% ------------------------------------------------------------------------------
\section{Synthese: Der Zielkonflikt bei Parser-Werkzeugen}
\label{sec:toolchain_synthesis}
% ------------------------------------------------------------------------------

Die Analyse verdeutlicht das klassische Spannungsfeld bei der Wahl von Parser-Werkzeugen in der Lehre und zeigt, wie Larkita es verschiebt.
Mit den früheren Werkzeugketten unterschieden sich die drei Fassungen nicht nur in der Sprache, sondern auch im Werkzeug: Python mit Lark bot maximale interaktive Leichtigkeit ohne statische Typsicherheit, OCaml mit Menhir herausragende Konfliktdiagnostik, statische Typsicherheit der semantischen Aktionen und den höchsten Durchsatz um den Preis externer Compiler-Prozesse, und F\# mit FsLexYacc professionelle Software-Engineering-Disziplin, die jedoch das spielerische Ausprobieren im Seminarraum erschwerte.

Mit Larkita und seinen Portierungen lesen alle drei Fassungen dieselbe Grammatik, erzeugen denselben Automaten, schreiben dieselbe Ausgabedatei und liefern dieselben Syntaxbäume und Fehlermeldungen.
Der Vergleich der Sprachen wird dadurch schärfer: Er betrifft nicht mehr die Werkzeuge, sondern nur noch die Darstellung und Verarbeitung der Syntaxbäume.
Python verarbeitet die heterogenen Tupel direkt; OCaml und F\# benötigen einen Summentyp, der ebenso wenig Struktur garantiert wie die Tupel, gewinnen ihre statische Sicherheit aber dort zurück, wo der Baum in getypte algebraische Datentypen umgewandelt wird.
Diese Umwandlung ersetzt die semantischen Aktionen von Menhir und FsYacc; sie steht nun im Notizbuch, sichtbar neben dem übrigen Code.
Den Preis zahlen die Portierungen beim Durchsatz, der für die Lehre ausreicht, aber deutlich unter dem generierter Parser liegt.
